\documentclass[12pt]{article}
\usepackage[utf8]{inputenc}
\usepackage[T1]{fontenc}
\usepackage{hyperref}
\title{Turaka Language Resources for K-12 Curriculum Development}
\author{Compiled by Claude Sonnet 5 (claude-sonnet-5) for Andr\'as Kornai}
\date{}
\begin{document}
\maketitle
\section{Introduction}
Turaka (ISO 639-3 trh) is a Trans-New Guinea language of the Dagan group, spoken in Milne Bay Province, Papua New Guinea. Speaker numbers differ between sources: 25 (cited 2000), 35 (reported 1981), and about 30 now. UNESCO classifies it as critically endangered. It is under pressure from larger related languages, especially Daga. Almost nothing exists in print or online.
\section{Primary Educational Resources}
I found none.
\section{Linguistic and Reference Materials}
\begin{itemize}
\item Overview of Turaka: \url{https://en.wikipedia.org/wiki/Turaka_language}
\item Basic wordlist in the ASJP database: \url{https://asjp.clld.org/languages/TURAKA}
\item The \emph{Turaka Preallocation Survey Report} (SIL-PNG, Ukarumpa, 1998) is cited in the Wikipedia article. I did not find a URL for it.
\end{itemize}
\section{Religious and Community Resources}
\begin{itemize}
\item Joshua Project people-group profile. It lists a population of 50 and notes that the group may have assimilated into larger groups: \url{https://joshuaproject.net/people_groups/15651}
\end{itemize}
\section{Digital and Audio Resources}
\begin{itemize}
\item Endangered Languages Project entry: \url{https://endangeredlanguages.com/lang/2695}
\end{itemize}
\section{Curriculum Development Context}
The language is not known to be taught in schools. With few speakers, all older adults, a curriculum is not realistic without urgent documentation first. The only educational option indicated is the larger neighbouring languages.
\end{document}